\documentclass[11pt]{article}
\usepackage[margin=1in]{geometry}
\usepackage{amsmath}

\title{A Transformer-Based Approach for Forecasting\\Emerging Research Trends in Scientific Literature}
\author{Hoang Tien Dat \\ FPT University}
\date{\today}

\begin{document}

\maketitle

% ============================================================
% DEMO STEP 1: Tab "Overlap Check" se tu trich abstract duoi day
% (khong can boi den) va so voi corpus -> ra danh sach bai trung.
% ============================================================
\begin{abstract}
The rapid growth of scientific publications makes it increasingly difficult
for researchers to identify emerging topics before they become mainstream.
In this paper, we propose a transformer-based deep learning framework that
forecasts emerging research trends by analyzing keyword co-occurrence
networks extracted from large-scale bibliographic databases. Our approach
combines a pretrained language model for semantic keyword embedding with a
temporal graph neural network that captures the evolution of topic
communities over time. We evaluate the framework on a corpus of publications
in machine learning and natural language processing, and show that it
predicts breakout keywords twelve months in advance with higher accuracy
than frequency-based baselines. These results suggest that semantic trend
forecasting can help researchers, students, and funding agencies position
their work earlier in the research life cycle.
\end{abstract}

\section{Introduction}
Keeping track of where a research field is heading is a challenge for both
newcomers and experienced scholars. Existing tools mostly rank topics by
raw publication counts, which lags behind the actual emergence of new ideas.
% DEMO STEP 2: dat con tro sau chu "attention mechanisms" roi bam "Cite"
% o tab Overlap Check -> \cite{...} duoc chen tai day va \bibitem tu dong
% xuat hien trong thebibliography cuoi file.
Our work builds on recent advances in attention mechanisms and graph
representation learning.

\section{Methodology}
Let $G_t = (V, E_t)$ denote the keyword co-occurrence graph at time step
$t$, where each node $v \in V$ is a keyword and each edge weight represents
the number of papers mentioning both keywords. We embed every keyword with
a pretrained encoder and update the temporal state with
\begin{equation}
  h_v^{(t)} = \mathrm{GRU}\!\left(h_v^{(t-1)},\;
  \sum_{u \in \mathcal{N}(v)} \alpha_{uv}\, W h_u^{(t-1)}\right),
\end{equation}
where $\alpha_{uv}$ is an attention coefficient over neighboring keywords.

\section{Experiments}
We collect publication metadata from an open bibliographic database and
split it chronologically: models observe five years of history and predict
which keywords will double their publication share within the next year.

\section{Conclusion}
We presented a semantic framework for research trend forecasting. Future
work includes modeling author collaboration networks and extending the
evaluation to additional scientific domains.

% ============================================================
% DEMO STEP 3: cac \bibitem do nut "Cite" sinh ra se duoc chen
% ngay truoc \end{thebibliography} ben duoi.
% ============================================================
\begin{thebibliography}{99}
\bibitem{vaswani2017} Ashish Vaswani, Noam Shazeer, Niki Parmar et al.
Attention Is All You Need. \textit{Advances in Neural Information
Processing Systems}, 2017.
\end{thebibliography}

\end{document}
